% Part: lambda-calculus
% Chapter: lambda-definability
% Section: introduction

\documentclass[../../../include/open-logic-section]{subfiles}

\begin{document}

\olfileid{lam}{ldf}{int}
\olsection{పరిచయం}

మొదట చూసినప్పుడు, లాంబ్డా కలనశాస్త్రం ప్రమేయాలను, వాటిని
ఇతర ప్రమేయాలకు ప్రయోగించడాన్ని సూచించే పదాలపై అత్యంత
అమూర్తమైన కలనశాస్త్రం మాత్రమే అనిపిస్తుంది. దాని
వాక్యనిర్మాణంలో అవి ప్రత్యేక రకమైన వస్తువులపై పనిచేసే
ప్రమేయాలనే సూచన లేదు; ముఖ్యంగా సహజ సంఖ్యల ప్రస్తావన
అందులో లేదు. ప్రయోగం, లాంబ్డా అమూర్తీకరణ అనే ప్రాథమిక
క్రియలు సహజ సంఖ్యలపై ప్రమేయాలకే కాక ఏ ప్రమేయానికైనా
వర్తిస్తాయి.

అయినప్పటికీ, కొంత చాతుర్యంతో లాంబ్డా కలనశాస్త్రంలో
అంకగణిత ప్రమేయాలను, అంటే సహజ సంఖ్యలపై ప్రమేయాలను,
నిర్వచించవచ్చు. అందుకోసం ప్రతి సహజ సంఖ్య~$n \in
\Nat$కు ప్రత్యేక $\lambd$-పదం~$\num n$ను నిర్వచిస్తాం;
అది $n$కు సంబంధించిన \emph{చర్చ్ సంఖ్యాంకం}.
(ఈ సంఖ్యాంకాలకు అలోన్జో చర్చ్ పేరు పెట్టారు.)

\begin{defn}
  $n \in \Nat$ అయితే, దానికి సంబంధించిన \emph{చర్చ్ సంఖ్యాంకం}
  $\num{n}$ అనేది $n$ను సూచిస్తుంది:
  \[
    \num{n} \ident \lambd[fx][f^n(x)]
  \]
  ఇక్కడ $f^n(x)$ అంటే $f$ను $x$కు $n$ సార్లు
  ప్రయోగించిన ఫలితం. ఉదాహరణకు, $\num{0}$ అనేది
  $\lambd[fx][x]$, $\num{3}$ అనేది
  $\lambd[fx][f(f(f\,x))]$.
\end{defn}

చర్చ్ సంఖ్యాంకం $\num n$ రెండు ఆర్గ్యుమెంట్లు $f$, $x$లను
స్వీకరించి $f^n(x)$ను తిరిగి ఇచ్చే ప్రమేయాన్ని సూచించే
లాంబ్డా పదంగా సంకేతీకరించబడింది. చర్చ్ సంఖ్యాంకాలు
స్పష్టంగానే నియత రూపంలో ఉన్నాయి.

సహజ సంఖ్యల లాంబ్డా ప్రాతినిధ్యం వాటితో గణించగలిగితేనే
ఉపయోగకరం. చర్చ్ సంఖ్యాంకాలతో గణించడం అంటే ఒక $\lambd$-
పదం~$F$ను అటువంటి సంఖ్యాంకానికి ప్రయోగించి, ఏర్పడిన
పదం~$F\, \num n$ను నియత రూపానికి తగ్గించడం. అది ఎప్పుడూ
నియత రూపానికి తగ్గి, ఆ నియత రూపం ఎప్పుడూ చర్చ్
సంఖ్యాంకం~$\num m$ అయితే, గణన ఫలితాన్ని సంఖ్య~$m$గా
భావించవచ్చు. అప్పుడు $F$ ఒక ప్రమేయం
$f\colon \Nat \to \Nat$ను నిర్వచిస్తుందని భావించవచ్చు:
$F\, \num n \red \num m$ అయితేనే $f(n)=m$.
చర్చ్--రోసర్ లక్షణం వల్ల నియత రూపాలు ఉంటే అవి అనన్యమైనవి.
కనుక $F\, \num n \red \num m$ అయితే, $F \, \num n$ తగ్గే
మరే ఇతర నియత రూపమూ, ముఖ్యంగా మరే ఇతర చర్చ్ సంఖ్యాంకమూ,
ఉండదు.

మరోవైపు, $f\colon \Nat \to \Nat$ అనే ప్రమేయం ఇస్తే,
$f$ను ఇదే విధంగా నిర్వచించే $F$ పదం ఉందా అని అడగవచ్చు.
ఉంటే, $F$ ఆ $f$ను \emph{\tetoken{లాంబ్డాతో నిర్వచిస్తుంది}{!!{lambda define}s}}
అనీ, $f$ \tetoken{లాంబ్డాతో నిర్వచించదగిన}{!!{lambda definable}}దనీ అంటాం.
దీన్ని బహుస్థానిక, పాక్షిక ప్రమేయాలకూ విస్తరించవచ్చు.

\begin{defn}
$f\colon \Nat^k \to \Nat$ అనుకుందాం. ప్రతి
$n_0$, \dots,~$n_{k-1}$కు $f(n_0, \dots, n_{k-1})$
నిర్వచితమైతే
\[
F \, \num{n_0} \, \num{n_1} \dots \num{n_{k-1}} \red \num{f(n_0, n_1, \dots,
  n_{k-1})}
\]
అవుతూ, అది నిర్వచితం కాకపోతే
$F \, \num{n_0} \, \num{n_1} \dots \num{n_{k-1}}$కు
నియత రూపం లేకపోతే, లాంబ్డా పదం~$F$ ప్రమేయం~$f$ను
\emph{\tetoken{లాంబ్డాతో నిర్వచిస్తుంది}{!!{lambda define}s}} అని అంటాం.
\end{defn}

అత్యంత సరళమైన ఉదాహరణ స్థిర ప్రమేయాలు. $C_k \ident
\lambd[x][\num{k}]$ అనే పదం, $c_k\colon \Nat \to
\Nat$ మరియు $c_k(n)=k$ అయిన స్థిర ప్రమేయాన్ని
\tetoken{లాంబ్డాతో నిర్వచిస్తుంది}{!!{lambda define}s}.
\sourcecorrection{OLTELAMLDFI-001}{మూలం స్థిర ప్రమేయాన్ని $c_k$గా పరిచయం చేసి, వెంటనే దాని విలువను $c(n)=k$ అని రాసింది. అదే ప్రమేయాన్ని సూచించడానికి తెలుగు వాక్యంలో $c_k(n)=k$ అని ప్రధానచిహ్నాన్ని పునరుద్ధరించాం; మిగిలిన తగ్గింపు గణన యథాతథం.}
ఏ $n$కైనా $C_k \, \num n \ident
(\lambd[x][\num{k}])\num n \redone \num{k}$.
తత్సమ ప్రమేయం $\lambd[x][x]$చే
\tetoken{లాంబ్డాతో నిర్వచించబడిన}{!!{lambda defined}}ది.
మరింత సంక్లిష్టమైన ప్రమేయాలను నిర్వచించడం కష్టమే;
తరచుగా చాలా చాతుర్యం అవసరం. కాబట్టి ప్రతి గణనీయ ప్రమేయం
\tetoken{లాంబ్డాతో నిర్వచించదగిన}{!!{lambda definable}}ది
అనేది ఆశ్చర్యకరంగా అనిపించవచ్చు. తిరుగు వాదన కూడా నిజమే:
ఏదైనా ప్రమేయం \tetoken{లాంబ్డాతో నిర్వచించదగిన}{!!{lambda definable}}దైతే,
అది గణనీయమైనది.

\end{document}
